让标注员给一段文本打一到十分，同一批文本发给两组人，平均分能差一分半；隔一周再发给同一组人，同一个人的打分也会漂将近一分。把任务换成``这两段哪个更好''，重测一致率能到八成以上。人对相对判断比对绝对判断稳定得多——这是可以反复观察到的经验规律，也是把训练信号建立在比较之上的全部理由。

问题随之而来：一堆两两比较，究竟能把``好''定到什么程度？答案不是``数据够多就能定''。有三处空缺是结构性的，数据量补不上。分数只由差值决定，整体平移在似然上完全不可见；比较图若分成互不相连的几团，团与团之间的关系没有任何数据约束；而把偏好压进一个实数分数，这个动作本身已经假设了偏好可传递，而真实的成对偏好可以有环。

本章把这三条限制先在最简单的 Bradley--Terry 模型上讲清楚，再看它们在奖励模型里以什么形式重新出现。换掉的只是参数化方式——把每个物品一个自由参数换成一个以输入为条件的函数——限制一条都没少，只是变得更难看见。由此可以读出几条具体判断：奖励的绝对尺度是未定义的量，任何依赖它的下游做法都在使用一个换个随机种子就会变的数；奖励模型学到的是标注者偏好的分布，标注指南、标注员构成与采样设计全都写在里面；而优化会把策略推出奖励模型有证据的区域，于是奖励继续涨、被奖励衡量的东西开始跌。

最后一节把视角从一个标注员扩到一群。把所有人的比较不加区分地加进同一个损失，这个求和动作看上去只是记账，实际上完整指定了一条偏好聚合规则。而 Arrow 的定理说明，满足帕累托一致、无关方案独立、无独裁的传递聚合规则不存在——所以任何可用的规则都已经在某处让步，``众包取多数''不是中立操作。

\subsection{三节怎么安排}

第一节建立 Bradley--Terry 模型，并把可辨识性、比较图连通性、传递性这三条限制逐一说清，每一条都指出它是代数结构、图结构还是关于世界的假设。第二节把这个似然写成奖励模型的训练损失，说明它继承了什么、丢掉了什么，并交代奖励过优化这个经验现象的结构来源。第三节从 Condorcet 悖论进入 Arrow 不可能定理，逐条对账三个前提各挡住什么失效模式，再回到损失函数是怎样悄悄选定一条聚合规则的。

\subsection{本章篇目}

\begin{itemize}
\tightlist
\item \hyperref[sec:17-Constrained-Guarantees/05-Learning-from-Comparisons/01]{``成对比较能定出什么''}；
\item \hyperref[sec:17-Constrained-Guarantees/05-Learning-from-Comparisons/02]{``从偏好到奖励模型''}；
\item \hyperref[sec:17-Constrained-Guarantees/05-Learning-from-Comparisons/03]{``把很多人的偏好合起来''}。
\end{itemize}

读这一章用得上三样前置：logistic 模型的形状，见\hyperref[sec:13-Machine-Learning/02-Linear-Models/03]{``从二分类到 Softmax 分类''}；参数辨识不出来是怎么被诊断的，见\hyperref[sec:11-Mathematical-Statistics/03-Sufficient-Statistics-and-Information/02]{``Score 与 Fisher 信息衡量局部可辨认性''}；图的连通分支意味着什么，见\hyperref[sec:04-Discrete-Mathematics/03-Graph-Theory/02]{``路径、环与连通性''}。第二节涉及的分布偏移可以对照\hyperref[ch:136]{``强化学习：模型未知时怎样逼近最优决策''}里的离策略部分读。
